% !TEX program = xelatex
% Chladni 板逆向设计系统：五个算法改进方向（含试错史回顾，通俗版）
% 编译：xelatex improvement_directions_zh.tex
\documentclass[11pt,a4paper]{article}

\usepackage[margin=2.5cm]{geometry}
\usepackage{amsmath,amssymb}
\usepackage{booktabs}
\usepackage{enumitem}
\usepackage{array}
\usepackage{xeCJK}
\setCJKmainfont{Droid Sans Fallback}[AutoFakeBold=2.5,AutoFakeSlant=0.15]
\setCJKsansfont{Droid Sans Fallback}[AutoFakeBold=2.5]
\setCJKmonofont{Droid Sans Fallback}
\usepackage[colorlinks=true,linkcolor=blue,urlcolor=blue]{hyperref}

\newcommand{\code}[1]{\texttt{#1}}
\linespread{1.18}
\setlength{\parskip}{0.35em}

\title{\textbf{Chladni 板逆向设计系统：五个算法改进方向\\
\large（含试错史回顾与勘误）}}
\author{Chladni Studio 项目组}
\date{2026 年 6 月}

\begin{document}
\maketitle

\begin{abstract}
我们的软件做的事情可以用一句话概括：用户画一个图案，程序反过来设计一块厚薄不均的板子和几个驱动频率，让沙子在振动的板上尽量摆出这个图案。经过 350 多次实验，目前的流水线已经能在可打印的碳纤维板（CF-PETG）上把 IC 字母目标做到"目标线上的沙密度是全板平均的 4.85 倍"，接近这套物理装置的极限。本文用尽量通俗的语言说明五个还能改进的方向：(1)~评分系统重构---让分数真正反映"看起来像不像"；(2)~几何趋势感知---让程序知道沙子偏在哪、该往哪挪；(3)~频率内循环---先给板子找最合适的频率，再评价板子好不好；(4)~随机重启---对结果不满意就换个起点重来；(5)~中心伪影修正---修复"仿真总在板中心堆一坨假沙"的建模错误。文章最后回顾了项目试错史中的几个关键转折，并勘正了其中流传的几处不准确的说法。
\end{abstract}

\section{背景：系统现在是怎么工作的}

整条流水线分四步。第一步，用户给一张 $256\times 256$ 的黑白目标图。第二步，一个我们自己写的"草稿物理引擎"（业内叫 surrogate，替代模型）快速迭代约 300 步，同时调整板子的 $15\times 15$ 厚度分布 $H$ 和 6 个候选驱动频率，让预测的沙图尽量接近目标。第三步，把这块板子交给精确但缓慢的 COMSOL 仿真，算出它的前 30 个共振模态，再从中挑出对目标贡献最大的几个频率（这一步叫 Phase~1），外加一个靠人工扫频发现的"幸运频率"（165~Hz）。第四步（Phase~2），拿 COMSOL 的精确结果反过来微调板子，小步快试：变好就继续，变差就退回，只跑一轮。

同样重要的是认清\textbf{物理边界}---有些"不像"不是算法的错：

\begin{itemize}[leftmargin=2em]
\item \textbf{中心单点激励}：激振器只顶在板中心一个点上，能有效推动的振型天然偏向"最对称的那一族"（群论里叫 A1 表示）。
\item \textbf{正方形对称}（D4 对称）：正方形板有 8 种把自己变回自己的旋转和翻转，板子和激励都对称时，振出来的图案也强烈倾向对称。非对称目标（比如字母 IC）因此先天吃亏。
\item \textbf{分辨率}：厚度只有 $15\times 15$ 格可调；材料各向异性强度（刚度比）受打印工艺限制，上限约 3。
\end{itemize}

在这些约束下，$4.85$ 倍富集已接近实用上限。下面五个方向不是要打破物理，而是：让评分更诚实、让搜索更聪明、让仿真更接近真实的沙。

\subsection*{名词速查}

\begin{center}
\small
\begin{tabular}{@{}p{0.26\textwidth}p{0.66\textwidth}@{}}
\toprule
术语 & 大白话解释 \\
\midrule
surrogate（替代模型） & 自写的快速近似物理引擎：算得糙但快，用来代替缓慢的 COMSOL 做大量迭代 \\
enrichment（富集度 $E$） & 目标线上的沙密度是全板平均的几倍。越大越好 \\
recall（覆盖率 $R$） & 目标线有多大比例真的落在"安静区"（沙会留下的低振幅区）里 \\
contrast（对比度 $C$） & 目标区沙密度与背景区沙密度之比 \\
节线 / 安静区 & 板上振幅接近零的地方，沙子最终聚集在这里 \\
D4 对称 / A1 振型 & 正方形的 8 种自我重合操作 / 在所有这些操作下都不变的最对称振型家族 \\
信任域 & "小步试探"策略：这步有改善就放大步长，变差就缩小步长退回 \\
离共振 / 幸运频率 & 不卡在共振峰上、却对目标图案贡献意外大的频率（如 165 Hz） \\
PCA 设计流形 & 从历史好设计中用主成分分析提炼出的几十个"有效改法"方向 \\
最优传输（OT） & 把当前沙图"搬"成目标图最少要花多少力气；搬运方案告诉你每堆沙该往哪挪 \\
\bottomrule
\end{tabular}
\end{center}

\section{方向一：评分系统重构}

\subsection{现状与问题}

程序怎么知道一块板"好不好"？靠打分。目前优化时用的损失函数是
\begin{equation}
\mathcal{L} \;=\; -\log E \;-\; \log C \;-\; 0.5\,\log R
\;\;\equiv\;\; -\log\!\bigl(E\cdot C\cdot R^{0.5}\bigr),
\end{equation}
本质上是把富集度、对比度、覆盖率\textbf{乘}在一起。排序时另用一个加权和：
\begin{equation}
S_{\mathrm{comp}} = 0.40\,\mathrm{clip}\bigl(\log(1{+}\max(E{-}1,0)),0,2\bigr) + 0.30\,R + 0.20\,\mathrm{clip}\bigl(\log(1{+}\max(C{-}1,0)),0,2\bigr) + 0.10\,D,
\end{equation}
其中 $D$ 衡量沙图的主方向是否和目标一致。

实测的问题是\textbf{分数高不等于看起来像}。最典型的例子：目标是一条对角线，某块板拿到了 $2.41$ 倍富集的好分数，但实际图案是个 X。原因不难理解：富集度只问"目标线上沙密不密"，不问"是不是整条线都有沙"---只要沙强烈集中在目标线的\textbf{一小段}上，分数照样高，哪怕图案整体完全不像。

\subsection{改进方案}

思路转换：富集度和对比度不再参与"谁更好"的排名，只设一个手动下限当\textbf{及格线}（低于及格线直接淘汰）；及格的候选之间，按"像不像"来排。

但"像不像"不能只用覆盖率 $R$。覆盖率有个盲区：它只问"目标线被盖住了多少"，不管目标\textbf{之外}有没有多余的杂线---而看起来不像的主要原因恰恰是杂线。所以排序键应当同时考虑两个方向："目标线被盖住多少"（recall）和"安静区里有多大比例真的落在目标上"（precision，查准率）。两者的标准组合就是 $F_\beta$ 分数：
\begin{equation}
S_{\mathrm{rank}} = F_\beta(\mathrm{valley},\,\mathrm{target})
\quad\text{s.t.}\quad E \ge \tau_E,\; C \ge \tau_C .
\end{equation}

三个实施细节：(a)~优化过程需要处处可导，"及格线"这种硬门槛不可导，改用铰链罚（离及格线越远罚越重，过线罚为零）；(b)~修一个已发现的小 bug：振幅全零的退化场会被误判为 $R=1.0$（满分），应判零分；(c)~及格线数值要按目标图案的粗细面积校准，不能全局一刀切。

\section{方向二：几何趋势感知评分}

\subsection{动机}

现在的三个指标都是"位置盲"的：它们能告诉你"有多少沙落在目标上"，却不能告诉你"沙偏在哪里、该往哪个方向挪"。两个分数完全相同的图案，几何上可以天差地别。这就像只告诉学生考了 60 分，却不告诉他错在哪---程序自然没有方向感。让算法"像人一样理解几何趋势"，数学上有现成的严格工具，而且\textbf{代码库里已经有两个半成品}：

\subsection{改进方案}

\begin{enumerate}[leftmargin=2em]
\item \textbf{距离场损失}（骨架见 \code{src/scoring/signed\_distance\_loss.py}）：给每粒沙记一笔"离目标线多远"的账，
\begin{equation}
\mathcal{L}_{\mathrm{SDF}} = \frac{\sum_{x} p(x)\, d_T(x)}{\sum_x p(x)},
\end{equation}
$p$ 是沙密度，$d_T$ 是到目标线的距离。沙离得越远罚越重，于是梯度天然把沙往目标方向"拉"。
\item \textbf{最优传输损失}（骨架见 \code{src/physics/sinkhorn\_loss.py}）：计算"把当前沙图搬成目标图"的最小搬运代价。它的搬运方案就是字面意义上的"前进方向"---哪堆沙该往哪挪、挪多远，全写在里面。
\item \textbf{多尺度比对}：在几个不同模糊程度的版本上同时打分---先看"大形状对不对"，再抠细节。这接近人眼的判断顺序，也让优化的路更平缓。
\item \textbf{拓扑体检}（只用于最终排序）：数一数图案有几个连通块、骨架有几个分叉。一条对角线和一个 X，富集度分不出来，分叉数一数就分出来了。
\end{enumerate}

要回应一个历史顾虑：项目早期试过 IoU / Dice 这类逐像素比对，失败后放弃了。但当时失败的原因是拿它们去比\textbf{离散的节线掩膜}（不可导、模态一切换就跳变）；现在的沙密度场是连续可导的，在它上面加距离场 / 最优传输损失，不踩当年的坑。分工上：新损失负责\textbf{指方向}（优化用），富集度 / 覆盖率继续负责\textbf{验收排序}，互补而非取代。

\section{方向三：频率内循环梯度逻辑}

\subsection{现状与问题}

目前优化器对厚度和频率\textbf{同时}做梯度下降。问题在于：板子的共振峰位置完全由厚度决定，厚度一改，所有共振峰都会移动，而频率变量只能慢吞吞地跟着挪。结果是：\textbf{一块本来很好的板，可能仅仅因为"当前这组频率不适合它"而得了低分，被优化器一票否决}。打个比方：评价一把吉他的好坏，却不允许先调弦---好琴也会被冤枉。

现有的缓解手段（前 60 步冻结频率、一次带 6 个频率、频率间隔惩罚等）都是治标。值得注意的是，流水线在 COMSOL 阶段其实已经做了正确的事---Phase~1 就是"先把这块板的所有模态都看一遍，再挑频率"---但它只在优化\textbf{结束后}做一次。

\subsection{改进方案}

把"先调弦再评价"搬进优化循环内部：
\begin{equation}
S(H) \;=\; \max_{f\,\in\,[f_{\min},\,f_{\max}]} s(H, f),
\end{equation}
即每一步都先给当前板子扫一遍全频段、找出它的最佳频率，在最佳频率处打分，再对厚度求梯度。数学上有 Danskin 定理背书：在最优频率处求出的梯度就是正确的梯度，理论上是干净的。

你们担心的"扫全频段太慢"有解：用\textbf{模态叠加法}。每步先把板子的振动模态解出来（一次特征分解），之后任意多个频率的响应都可以用这些模态拼出来，几乎不要钱---可能比现在逐个频率直接求解还快。两个细节：(a)~厚度变化时最佳频率会突然跳变（两个模态换位），硬取 max 在跳变点不连续，用"软 max"平滑即可；(b)~你们提出的"步子要小，让新板的最佳频率和旧板差不多"的直觉完全正确---这就是信任域思想，可以和 Phase~2 现成的"变好加步、变差退回"逻辑合并。另外，"从平板出发"已经是现状（厚度默认初始化为均匀板），"先判断目标可不可能做出来"的预检也已存在（\code{target\_realizability.py}），都可直接复用。

\section{方向四：结果评估与随机重启}

\subsection{现状与问题}

梯度下降是"下山"算法：走到一个坑底就停，不知道隔壁是不是有更深的坑（局部最优问题）。项目曾经建过"多起点"基础设施（同时从 N 个起点下山，取最好的），但后来 $\theta$ 优化被撤销、随机性消失，现在的优化器是\textbf{确定性}的---同一个起点跑一百遍结果一模一样，多起点已名存实亡（代码里已标注为空操作）。

\subsection{改进方案}

\begin{enumerate}[leftmargin=2em]
\item \textbf{先恢复随机性}：重新引入随机的厚度初始化（或对均匀板初值加受控扰动）。没有随机性，"重启"就是空话。
\item \textbf{换个聪明的空间重启}：与其在 225 个厚度格子上瞎撒点，不如在"PCA 设计流形"里采样---那是从历史好设计中提炼出的几十个有效方向。"离原起点足够远"用拉丁超立方等标准采样方法保证。
\item \textbf{及格线触发重启，但及格线不能拍脑袋}：IC 目标的 4.85 倍已接近物理上限，如果把及格线定得比物理上限还高，程序会永远重启下去。及格线应从"这个目标理论上最多能做到多少"的可实现性分析推导，按目标自适应，并设重启次数上限。
\item \textbf{复用现成的防退化合约}：已有的"新答案必须不差于老答案、换冠军需要超过阈值的提升"逻辑可以直接套用。
\end{enumerate}

这个方向要排在方向一、二之后---重启"靠什么判断不满意"和"靠什么比较好坏"，都依赖一套更可信的评分。

\section{方向五：中心激励伪影修正}

\subsection{问题确认}

实验中观察到一个仿真和现实的顽固矛盾：\textbf{现实里激励点（板中心）几乎永远不积沙---它是全板震得最凶的地方，沙最先被弹飞；而仿真输出每次都在中心堆一坨沙。}

检查代码后确认这是一个系统性的建模错误，机制如下。仿真采用"相对坐标系"描述基础激励：把中心夹持点的位移直接固定为零（\code{orthotropic\_plate.py} 中中心自由度被移除、\code{expand\_to\_grid} 把中心格填 0），等价于站在激振器上看板子。在这个视角里，中心的\textbf{相对}位移恒为零。而撒粉模型
\begin{equation}
p(x) = \exp\!\Bigl(-\bigl(|u(x)|/\sigma\bigr)^2\Bigr)
\end{equation}
的规则是"哪里振幅小，沙就堆哪里"---于是中心被误判为最安静的地方，堆出一坨\textbf{幻影沙}。现实中沙听的不是"相对位移"，而是\textbf{绝对加速度}：加速度低于临界值（约重力加速度 $g$ 的量级）沙才留得住，而中心的绝对加速度恰恰全板最大。

\subsection{改进方案}

\begin{enumerate}[leftmargin=2em]
\item \textbf{快修（一行代码级）}：打分时把中心夹持区直接挖掉不计分。注意 surrogate 和 COMSOL 评分（\code{score\_comsol\_response.py}）两边都要挖。
\item \textbf{正修（物理正确）}：打分改用绝对位移。谐波激励下基础位移幅值是 $a_{\mathrm{base}}/\omega^2$，
\begin{equation}
u_{\mathrm{abs}}(x) = u_{\mathrm{rel}}(x) + a_{\mathrm{base}}/\omega^{2},
\end{equation}
加上这一项后中心不再是零，幻影沙自动消失。
\item \textbf{撒粉模型升级}：把"相对位移小就积沙"换成"绝对加速度低于阈值才积沙"的软开关，如 $p \propto \sigma\!\bigl(\beta\,(a_{\mathrm{crit}}-\omega^2|u_{\mathrm{abs}}|)\bigr)$，这才是真实的沙物理。
\end{enumerate}

\subsection{连带后果}

幻影沙不是局部小瑕疵：它混进了沙的总量（富集度的分母）、扰动了"安静区"的判定阈值、压低了对比度---\textbf{所有历史分数都被它污染过}。修掉之后，4.85 倍这个基准一定会变（目标不含中心时大概率变好），幸运频率 165~Hz 和 Phase~1 的选模结果也可能要重新扫。因此这个方向应当\textbf{最先做}：它最便宜、全局受益；反过来，如果不先修它，方向二的最优传输损失会一本正经地把这坨假沙"搬运"到目标上---方向都给错了。

\section{试错史回顾与勘误}

项目组整理的试错记录（trails）保存了很多有价值的判断，但其中几处"为什么这么改"的解释与代码库实际记录的证据有出入。这一节把试错史和前五个方向接到一起，顺带勘正。

\subsection{材料之路：从"均匀"到"各向异性"，再到撤销 $\theta$}

试错记录的故事线是对的：最早用 SLA 打印追求材料均匀，结果发现\textbf{即使厚度差异做得很大，图案仍顽固地保持对称、和均匀板差别不大}；于是改用 FDM 打印各向异性材料；先尝试让每个小格子有不同的纤维取向（$\theta$ 场），后来发现效果不好，退回"全板同一取向"。

这条路线背后的物理可以说得更清楚。纯靠厚度破坏对称之所以弱，是因为弯曲刚度正比于厚度的\textbf{三次方}，但可打印的厚度范围有限，撬不动多少；而统一取向的各向异性直接把正方形的 8 重对称（D4）降到 4 重（D2）---所以图案能偏离均匀板，却仍保留镜面对称。这正好解释了"不同于均匀板、但还是挺对称"的观察。

需要勘正的是\textbf{撤销 $\theta$ 的原因}。试错记录写的是"自由度太多，程序混乱，计算变慢"。代码库记录的真实机制（W10 模块文档，引 2026-05 的对照实验）不同：W10 实际工作在\textbf{离共振}频段，在那里惯性项比刚度项大约 5 个数量级（$\omega^2 M \gg K$），而 $\theta$ 只通过刚度项进入方程---所以 $\theta$ 的\textbf{梯度近乎为零}。它不是"多到让程序困惑"，而是"物理上几乎不起作用"的寄生变量：优化器推不动它，它就随机漂移，漂移的 $\theta$ 又把下游 COMSOL 的各向异性模态搅得更糟。10 次对照实验（5 个目标 $\times$ 2 种材料）证实：各向同性 PLA 在 4/6 个目标上反而\textbf{赢过}带"优化 $\theta$"的碳纤维板。这个版本的因果链更准确，也更有说服力。

顺带一个文档问题：\code{docs/final\_results.zh-CN.md} 仍写着"关键洞察：$\theta\neq 0$ 打破 D4 对称"，与现行代码（$\theta$ 恒为零）矛盾，应同步修订。

\subsection{评分之路："只在中间堆沙"的再归因}

试错记录：最早的评分是"沙靠近目标加分、远离扣分"，结果程序总是产出"沙只堆在板中间"的设计；当时归因为"物理限制---方程不存在和目标一样的解"，于是改用 enrichment / recall / contrast 三指标。

这个归因只对了一半。结合方向五的发现，\textbf{早期"中间堆沙"很可能有相当一部分是中心幻影沙伪影}---建模错误把中心标成了最适合积沙的地方，程序自然把"产量"都堆到那里。物理限制是真实存在的，但它的表现是"图案偏对称、偏简单"，而不是"沙都挤到中心"。建议试错记录补上这层再归因，它把两次独立的困惑（早期评分怪象、现在的中心沙堆）用同一个根因串起来了。

另外，三指标的前传值得写进记录：在它们之前还有一代 IoU / Dice 逐像素评分（试错史 Phase~A--D），失败原因是离散不可导、模态一切换分数就跳变。是"撒粉物理"这个建模转换（沙堆在振幅低谷）救活了评分系统。这段历史也是方向二敢于重新引入"几何比对"的依据---这次是在连续可导的沙密度场上做，不是在离散掩膜上。

\subsection{surrogate 到底是什么}

试错记录把 surrogate 描述为"PyTorch 里的一个微分方程求解工具"。更准确的说法：它\textbf{不是现成工具}，而是我们自己写的板振动求解器（把 Kirchhoff--Love 板方程离散到粗网格上），选择用 PyTorch 实现的关键原因不是快，而是 PyTorch 的\textbf{自动求导}---它让"沙图对厚度的导数"可以直接算出来，整个梯度优化路线才成立。"快"来自三件事：网格粗（$25\times 25$）、放弃高精度、不追求和 COMSOL 一致，只求方向正确。COMSOL 仍是最终裁判，surrogate 只是先用草稿纸把方案摸出来。

\subsection{离共振的决定与它的连锁反应}

试错记录里"跳出特征模态"（不局限于共振频率，任何能出稳定图案的频率都可以用）是一个被金属板实验支持的好决定，幸运频率 165~Hz 就是它的直接产物。但值得指出一个当时没人意识到的连锁反应：\textbf{正是"转入离共振频段"这个决定，让 $\theta$ 变成了寄生变量}（离共振处 $\omega^2 M \gg K$，见 7.1 节）。两条看似独立的试错记录其实是一条因果链：先选择了离共振，后发现 $\theta$ 无用---后者是前者的必然结果。把这条链写明，能让试错史从"流水账"变成"有解释力的叙事"。

\section{实施顺序与预期收益}

\begin{table}[h]
\centering
\begin{tabular}{@{}clll@{}}
\toprule
阶段 & 方向 & 工作量 & 预期收益 \\
\midrule
0 & 五：中心伪影修正（先挖掉，再换绝对响应） & 小 & 全部指标去偏，基准重校 \\
1 & 一：评分重构（及格线 + 像不像排序） & 小--中 & 分数与肉眼判断对齐 \\
1 & 二：几何损失（距离场 / 最优传输 + 多尺度） & 中 & 程序知道沙该往哪挪 \\
2 & 三：频率内循环（先调弦，再评琴） & 中--大 & 不再冤枉好板子 \\
3 & 四：随机重启（聪明的换起点） & 中 & 跳出局部最优 \\
\bottomrule
\end{tabular}
\caption{实施路线图。阶段 0 是前置条件：不先清除幻影沙，新评分和几何损失都会在被污染的数据上工作。}
\end{table}

\section{结语}

五个方向针对同一条主线的三个薄弱环节：\textbf{评分的诚实度}（方向一、二、五）、\textbf{搜索的效率}（方向三、四）、\textbf{仿真的保真度}（方向五）。它们都不指望突破"单点激励 + 正方形对称"设定的物理天花板，而是保证算法在物理允许的范围内不浪费任何机会。试错史的回顾还给了一个方法论提醒：\textbf{把"为什么改"的真实证据留在记录里}（如 $\theta$ 的梯度消失、中心沙的建模根因），比只记"改了什么"重要得多---本项目几次最有价值的进展，都来自对旧结论的重新归因。推进次序上，先做便宜且全局受益的修正（中心伪影、评分键），再动优化器的结构（频率内循环、重启策略），风险最低。

\end{document}
